\phantomsection\label{fig:vol04:northern-song:early-state-capacity-transport}
\begin{center}
\begin{tikzpicture}[x=.88cm,y=.88cm,font=\small]
  \fill[paper] (0,0) rectangle (11.8,6.35);

  \fill[source!16] (.55,3.7) -- (3.55,3.5) -- (3.75,5.85) -- (.8,5.95) -- cycle;
  \node[align=center] at (2.12,4.8) {河东、河北\\边防与军需压力};
  \fill[timeline!21] (.7,.65) -- (4.0,.75) -- (3.7,2.9) -- (.5,3.0) -- cycle;
  \node[align=center] at (2.2,1.78) {四川\\山地道路、粮赋\\与接收后的治理};
  \fill[dynasty!19] (8.3,.65) -- (11.25,.8) -- (11.15,3.2) -- (8.1,2.95) -- cycle;
  \node[align=center] at (9.7,1.8) {江南\\水网、税粮、手工业\\与渡江战争};
  \fill[timeline!18] (8.1,3.65) -- (11.25,3.55) -- (11.3,5.9) -- (8.35,5.8) -- cycle;
  \node[align=center] at (9.85,4.85) {岭南\\港口、盐业与\\边地军政};

  \node[draw=dynasty!85,fill=dynasty!15,rounded corners=3pt,align=center,inner sep=5pt] (tokyo)
    at (5.85,3.15) {东京开封\\官署、仓储、禁军\\与文书调度};
  \node[draw=source!70,fill=paper,rounded corners=2pt,align=center,inner sep=3pt] (officials)
    at (5.85,5.18) {任命、轮调、文官监督\\与地方接收};
  \draw[<->,thick,color=source] (tokyo.north) -- node[right,font=\scriptsize]{调度、考核} (officials.south);

  \draw[->,ultra thick,color=source!85] (3.6,4.55) -- node[above,sloped,font=\scriptsize]{兵员、粮饷、信息} (5.15,3.65);
  \draw[->,ultra thick,color=timeline!90] (3.55,2.05) -- node[below,sloped,font=\scriptsize]{道路、赋役与接收} (5.12,2.75);
  \draw[->,ultra thick,color=dynasty!88] (8.05,1.95) -- node[below,sloped,font=\scriptsize]{水路、税粮与军需} (6.55,2.76);
  \draw[->,ultra thick,color=timeline!85] (8.05,4.65) -- node[above,sloped,font=\scriptsize]{港口、盐业与政务} (6.55,3.62);

  \node[text=source,align=center,font=\scriptsize] at (5.86,.75)
    {统一与守边依赖区域资源进入中枢；实际征收与运输并非单向无阻};
\end{tikzpicture}

\smallskip
\small\textit{图\quad 北宋初年区域资源、运输与国家能力关系示意：据官修纪传《宋史》卷1、卷2〈太祖本纪〉及后蜀、南汉、南唐世家，南宋编年汇集《续资治通鉴长编》乾德、开宝诸年条，后世辑存的制度汇编《宋会要辑稿·方域》，并参照开封北宋东京城遗址的水系考古资料，说明东京、边防、四川、江南和岭南之间的资源接收与调度关系。它不表示税额、运量、行政执行速度或某一地区必然服从中枢；图中区域边界、水道和资源、军需、信息路线均为约略示意。}
\end{center}
